\documentclass[../../../include/open-logic-section]{subfiles}

\begin{document}

\olfileid{sth}{spine}{valpha}

\olsection{टप्प्यांची $V_\alpha$ या रूपात व्याख्या}

\olref[sth][ordinals][]{chap} मध्ये आपण सुक्रमणे आणि (फोन न्यूमन यांच्या) क्रमसूचक संख्या परिभाषित केल्या. या प्रकरणात त्यांचा वापर करून संचांच्या पदानुक्रमाचे \emph{स्वतःचे} वर्णन करू. यासाठी \olref[sth][ordinals][opps]{sec} मधील उत्तरवर्ती आणि सीमा क्रमसूचक संख्यांच्या संकल्पना आठवू. पुढील व्याख्येत त्या वापरल्या आहेत:

\begin{defn}\ollabel{defValphas}
	\begin{align*}
	V_\emptyset &\defis \emptyset\\
	V_{\ordsucc{\alpha}} &\defis \Pow{V_\alpha} & &
	\text{प्रत्येक क्रमसूचक संख्या }\alpha\\
	V_{\alpha} &\defis \bigcup_{\gamma < \alpha} V_\gamma & &
	\text{जेव्हा }\alpha\text{ ही सीमा क्रमसूचक संख्या असते}
\end{align*}
\end{defn}
\noindent
ही क्रमसूचक संख्यांवरील \emph{अनंतपार पुनरावर्तनाने} केलेली व्याख्या ठरेल. तिची तुलना नैसर्गिक संख्यांवरील फलनांच्या पुनरावर्ती व्याख्यांशी करावी.\footnote{\olref[sfr][infinite][dedekind]{sec} मधील बेरीज, गुणाकार आणि घातांकनाच्या व्याख्या पाहा.} नैसर्गिक संख्यांच्या बाबतीत आधारप्रकरण आणि उत्तरवर्ती प्रकरणे परिभाषित करतो; पण क्रमसूचक संख्यांच्या बाबतीत \emph{सीमा} प्रकरणांतील वर्तनही सांगावे लागते.

$V_\alpha$ ची ही व्याख्या एक महत्त्वाचा टप्पा ठरेल. संचांची रचना \emph{टप्प्याटप्प्याने} होते, या कल्पनेने आपण संचांच्या पदानुक्रमाचे अनौपचारिक समर्थन केले. $V_\alpha$ हे प्रत्यक्षात तेच टप्पे आहेत. पण येथे महत्त्वाची बाब म्हणजे टप्प्यांचे हे \emph{आंतरिक} वर्णन आहे. उपपत्तीचे एखादे संभाव्य \emph{प्रतिमान} सुचवण्याऐवजी आपण टप्पे आपल्या संच उपपत्तीच्या \emph{आतच} परिभाषित करू.

\end{document}
